% 三种接口的平行对照，每条路径只表达对象交接。
\begin{tikzpicture}[
  font=\figurefont\footnotesize,
  rsbox/.style={rectangle,draw=BookRule,fill=BookPaper,
    align=center,text=BookInk,text width=24mm,minimum height=11mm,inner sep=3pt},
  rsflow/.style={-{Stealth[length=1.8mm]},draw=BookMuted,line width=.8pt},
  rshead/.style={font=\figurefont\bfseries\footnotesize,text=BookInk,anchor=west}
]
  \node[rshead] at (-1.3,0.9) {目录检索与候选评分};
  \node[rsbox] (r1) at (0,0) {目录访问};
  \node[rsbox] (r2) at (3.7,0) {候选集合\\逐候选评分};
  \node[rsbox] (r3) at (7.4,0) {列表组织\\核验与展示};
  \draw[rsflow] (r1.east) -- (r2.west);
  \draw[rsflow] (r2.east) -- (r3.west);
  \node[rshead] at (-1.3,-1.2) {物品标识生成};
  \node[rsbox] (s1) at (0,-2.1) {历史与条件\\生成物品代码};
  \node[rsbox] (s2) at (3.7,-2.1) {目录映射\\候选集合};
  \node[rsbox] (s3) at (7.4,-2.1) {评分及列表\\核验与展示};
  \draw[rsflow] (s1.east) -- (s2.west);
  \draw[rsflow] (s2.east) -- (s3.west);
  \node[rshead] at (-1.3,-3.3) {目录列表生成};
  \node[rsbox] (l1) at (0,-4.2) {历史与条件\\生成多物品代码};
  \node[rsbox] (l2) at (3.7,-4.2) {目录映射\\候选列表};
  \node[rsbox] (l3) at (7.4,-4.2) {列表核验\\执行展示};
  \draw[rsflow] (l1.east) -- (l2.west);
  \draw[rsflow] (l2.east) -- (l3.west);
\end{tikzpicture}
